% Shared TikZ helpers for the datapath figures (Harris & Harris style, ch. 7.5).
%   32-bit buses: thick black; narrow fields (Rs/Rd): thin black;
%   control / hazard signals: thin blue; additions of the dual-issue design: green.
\usetikzlibrary{calc}
\usetikzlibrary{circuits.logic.US}

\definecolor{ctl}{RGB}{11,111,184}    % control unit / hazard unit (blue, as in the book)
\definecolor{new}{RGB}{21,128,61}     % elements added for the dual-issue design
\definecolor{hzd}{RGB}{217,95,2}     % hazard unit and its wires (orange)
\definecolor{bru}{RGB}{124,58,237}   % shared branch unit (purple)

\tikzset{
  w/.style   = {line width=1.1pt},               % 32-bit datapath wire
  nw/.style  = {line width=0.45pt},              % narrow datapath field (5 bits)
  c/.style   = {line width=0.45pt, ctl},         % control wire
  n/.style   = {line width=0.45pt, new},         % dual-issue control wire
  h/.style   = {line width=0.45pt, hzd},         % hazard unit wire
  blk/.style = {line width=1.1pt, fill=white},
  lab/.style = {font=\fontsize{6.5}{6.5}\selectfont, inner sep=1pt},
  bit/.style = {font=\fontsize{5}{5}\selectfont, inner sep=0.6pt},
  mxl/.style = {font=\fontsize{4.3}{4.3}\selectfont, inner sep=0.4pt},
  ttl/.style = {font=\fontsize{7.5}{7.5}\selectfont\bfseries, inner sep=1pt, align=center},
}

% wire junction (T joints)
\newcommand{\dotat}[1]{\fill #1 circle (0.045);}

% \Mux[labels]{name}{x}{yc}{n}{d}{s}
%   vertical n:1 multiplexer pointing right, inputs spaced d.
%   s = +1: input 0 is the lowest one; s = -1: input 0 is the highest one.
%   labels: optional comma list written inside (default 0..n-1).
%   anchors: name-0 .. name-(n-1), name-o (output),
%            name-t / name-b (select, top / bottom),
%            name-so / name-si (select on the side of the last / first input)
\newcommand{\muxhw}{0.2}
\newcommand{\Mux}[7][]{%
  \pgfmathsetmacro\mh{(#5*#6+0.2)/2}%
  \draw[blk] (#3-\muxhw,#4-\mh) -- (#3-\muxhw,#4+\mh) -- (#3+\muxhw,#4+\mh-0.13)
             -- (#3+\muxhw,#4-\mh+0.13) -- cycle;
  \pgfmathtruncatemacro\mlast{#5-1}%
  \foreach \i in {0,...,\mlast} {
    \pgfmathsetmacro\my{#4 + #7*(\i-(#5-1)/2)*#6}
    \coordinate (#2-\i) at (#3-\muxhw,\my);
  }
  \def\mlabs{#1}%
  \ifx\mlabs\empty
    \foreach \i in {0,...,\mlast} {
      \node[mxl, anchor=west] at ($(#2-\i)+(0.02,0)$) {\i};
    }
  \else
    \foreach \lb [count=\i from 0] in {#1} {
      \node[mxl, anchor=west] at ($(#2-\i)+(0.02,0)$) {\lb};
    }
  \fi
  \coordinate (#2-o) at (#3+\muxhw,#4);
  \coordinate (#2-t) at (#3,#4+\mh-0.065);
  \coordinate (#2-b) at (#3,#4-\mh+0.065);
  \coordinate (#2-so) at (#3,#4+#7*\mh-#7*0.065);
  \coordinate (#2-si) at (#3,#4-#7*\mh+#7*0.065);
}

% \Alu{name}{x}{ylo}{yhi}{text}: ALU / adder shape (chevron with notch),
%   output name-o; text is written horizontally (use \AluV for a rotated "ALU")
\newcommand{\AluShape}[4]{%
  \pgfmathsetmacro\ac{(#3+#4)/2}\pgfmathsetmacro\ah{#4-#3}%
  \draw[blk] (#2,#3) -- (#2+0.9,#3+0.25*\ah) -- (#2+0.9,#4-0.25*\ah) -- (#2,#4)
             -- (#2,\ac+0.12*\ah) -- (#2+0.2,\ac) -- (#2,\ac-0.12*\ah) -- cycle;
  \coordinate (#1-o) at (#2+0.9,\ac);
}
\newcommand{\Alu}[5]{%
  \AluShape{#1}{#2}{#3}{#4}%
  \node[lab, font=\fontsize{7}{7}\selectfont\bfseries] at (#2+0.58,\ac) {#5};
}
\newcommand{\AluV}[5]{%
  \AluShape{#1}{#2}{#3}{#4}%
  \node[lab, font=\fontsize{8}{8}\selectfont, rotate=90] at (#2+0.6,\ac) {#5};
}

% \PipeReg{x}{yin}{ymid}{yout}{note}: pipeline register of one lane.
%   yin = edge next to the middle of the figure, yout = outer edge (clock side),
%   ymid = boundary between the datapath part (black) and the control part (blue).
%   As in the book the clock is always on the top edge (labelled CLK only when the
%   top edge is the outer one) and the note (e.g. CLR) is written at the bottom edge,
%   where the hazard unit signal comes in.
\newcommand{\PipeReg}[5]{%
  \draw[blk] (#1-0.175,#2) rectangle (#1+0.175,#3);
  \draw[line width=1.1pt, ctl, fill=white] (#1-0.175,#3) rectangle (#1+0.175,#4);
  \pgfmathsetmacro\pt{max(#2,#4)}\pgfmathsetmacro\pb{min(#2,#4)}%
  \draw[line width=0.5pt] (#1-0.12,\pt) -- (#1,\pt-0.15) -- (#1+0.12,\pt);
  \ifdim #4pt>#2pt \node[lab, anchor=south] at (#1,\pt+0.05) {CLK};\fi
  \node[bit, rotate=90, anchor=center] at (#1,\pb+0.3) {#5};
}

% \Bar{x}{ylo}{yhi}: plain register (clock on the top edge)
\newcommand{\Bar}[3]{%
  \pgfmathsetmacro\bt{max(#2,#3)}%
  \draw[blk] (#1-0.175,#2) rectangle (#1+0.175,#3);
  \draw[line width=0.5pt] (#1-0.12,\bt) -- (#1,\bt-0.15) -- (#1+0.12,\bt);
}
